\section{Masked Self Attention}

\subsection{为什么需要 Masking}

在语言建模等自回归任务中，生成第 $t$ 个 token 时\textbf{不应该看到}位置 $t+1, t+2, \ldots$ 的信息——否则就是``作弊''。但标准的 Self Attention 允许每个位置看到所有其他位置。

\begin{figure}[H]
\centering
\includegraphics[width=0.95\textwidth]{slides-images/slide-037.jpg}
\caption{Masked Self Attention：通过将未来位置的对齐分数设为 $-\infty$，在 softmax 后这些位置的注意力权重变为 0\protect\footnotemark}
\end{figure}
\footnotetext{Slides 第38页。}

\subsection{实现方式}

在计算完对齐分数 $E = \frac{QK^T}{\sqrt{d_q}}$ 之后：

\begin{enumerate}
    \item 构造一个\textbf{因果掩码}（causal mask）矩阵 $M$，其中位置 $(i,j)$：
    \begin{itemize}
        \item 若 $j \leq i$（允许看到），$M_{ij} = 0$
        \item 若 $j > i$（不允许看到），$M_{ij} = -\infty$
    \end{itemize}
    \item 将掩码加到对齐分数上：$E' = E + M$
    \item 执行 softmax：因为 $e^{-\infty} = 0$，被掩码的位置注意力权重为 0
\end{enumerate}

\begin{importantbox}{Masked Self Attention 用于语言建模}
所有现代的自回归语言模型（GPT 系列、LLaMA、Claude 等）都使用 Masked Self Attention（也称为 Causal Attention）。它确保模型在训练时只能依赖已经生成的 token，与推理时的自回归生成过程保持一致。

这也是所谓``decoder-only transformer''（仅解码器 Transformer）的核心——不需要单独的编码器和解码器，只需一个带因果掩码的 Self Attention 层栈。
\end{importantbox}

\subsection{本章小结}

Masked Self Attention 通过在注意力分数中注入 $-\infty$ 来阻止信息的``前向泄露''，使得 Self Attention 可以用于自回归生成任务。这是现代语言模型的基础构件。

%% ========== Multi-Head Attention ==========

